\chapter{Cơ sở lý thuyết}
\label{chap:background}

Chương này trình bày nền tảng khái niệm và các công nghệ được sử dụng trong
báo cáo. Phần đầu tóm tắt đặc trưng Big Data và lý do dữ liệu cảm biến IoT
là một trường hợp tiêu biểu. Phần thứ hai xây dựng mô hình xử lý event-time
có nhận thức veracity (veracity-aware event-time) làm hệ quy chiếu cho toàn
bộ báo cáo. Phần thứ ba mô tả chức năng, cơ chế hoạt động và lý do lựa chọn
của từng thành phần công nghệ trong kiến trúc
Chương~\ref{chap:implementation}, theo chiều luồng dữ liệu: lớp thiết bị,
lớp đệm, lớp xử lý, lớp lưu trữ --- phục vụ và lớp điều phối batch. Phần
cuối tổng hợp các nghiên cứu liên quan --- gồm các thực hành benchmark cho
hệ thống event-streaming và các mô hình giám sát chất lượng dữ liệu
stream-first --- kèm bảng so sánh khoảng trống nghiên cứu.

Một quy ước đọc quan trọng xuyên suốt phần mô hình: mọi định nghĩa, metric
và bất biến được gắn một trong ba trạng thái bằng chứng --- \textbf{baseline
đã triển khai} (cơ chế tương ứng đang chạy trong
Chương~\ref{chap:implementation}), \textbf{có thể suy ra từ dữ liệu đã ghi}
(tính được từ log hoặc dòng đã lưu mà không cần thêm cơ chế mới), hoặc
\textbf{extension chưa triển khai} (định nghĩa ở mức khái niệm, chưa có mã
chạy và chưa có số đo). Các định nghĩa extension không được dùng để mô tả
hành vi của baseline.

\section{Tổng quan Big Data và đặc thù dữ liệu cảm biến IoT}

Big Data thường được mô tả qua các đặc trưng Volume (khối lượng lớn),
Velocity (tốc độ sinh dữ liệu cao), Variety (đa dạng định dạng),
Veracity (độ tin cậy không đồng nhất) và Value (giá trị khai thác được).
Trong bối cảnh IoT sensor streaming, các đặc trưng thường bị ảnh hưởng mạnh
nhất là \textbf{Volume}, \textbf{Velocity} và \textbf{Veracity}: một cụm cảm
biến công nghiệp hoặc smart-city có thể sinh từ hàng triệu đến hàng tỷ event
mỗi ngày với
tần suất giây đến mili-giây, dung sai xử lý cuối cùng thường chỉ tính bằng
giây, và một tỷ lệ bản tin không hoàn hảo (muộn, trùng, ngoài ngưỡng vật
lý) luôn tồn tại trong luồng dữ liệu~\cite{iot_big_data_survey_2022}.

Các khảo sát về hệ thống phân tích IoT Big Data cũng chỉ ra rằng đường cong
tăng trưởng dữ liệu IoT kéo theo ba nhóm thách thức: (i) kiến trúc ingestion
chịu tải cao, (ii) lưu trữ có hỗ trợ ghi liên tục, và (iii) mô hình xử lý có
thể phục hồi sau sự cố~\cite{iot_big_data_survey_2022}. Đây chính là ba trụ
cột mà hệ thống được thiết kế trong báo cáo này nhắm tới; chiều Veracity
được bổ sung xuyên suốt như một trục thiết kế thay vì một bước hậu xử lý.

Khác với bài toán phân tích log/microservices truyền thống, dữ liệu cảm biến
có các tính chất khiến việc xử lý phải được thiết kế cẩn trọng ngay từ đầu:

\begin{itemize}
  \item \textbf{Tốc độ cao liên tục}: một dây chuyền với vài nghìn sensor,
        mỗi sensor đo 1--10 Hz, tạo ra hàng chục nghìn event/s;
  \item \textbf{Tính không hoàn hảo của timestamp}: event sinh ra tại cảm biến
        nhưng được nền tảng quan sát muộn hơn một lượng không hằng; event có
        thể đến sai thứ tự sinh, và chính timestamp cũng có thể lệch do clock
        cục bộ (xem Mục~\ref{subsec:lateness});
  \item \textbf{Trùng lặp}: do cơ chế retry hoặc idempotency yếu ở tầng gateway;
  \item \textbf{Giá trị giảm theo thời gian}: một cảnh báo nhiệt độ vượt ngưỡng
        gần như vô giá trị nếu đến tay người vận hành sau vài phút;
  \item \textbf{Tỷ lệ bản tin hỏng không đồng nhất}: noise, lỗi cảm biến,
        dropout kết nối tạo ra một tỷ lệ bản tin cần xử lý riêng, và tỷ lệ
        này khác nhau theo từng loại chỉ số đo.
\end{itemize}

Những tính chất này biến bài toán cảm biến IoT thành một use-case điển hình
của stream processing, nơi các khái niệm như event-time, idempotency và
fault tolerance trở thành yêu cầu cốt lõi
chứ không phải tuỳ chọn~\cite{iot_stream_latency_2021}.

\section{Mô hình hình thức: xử lý event-time có nhận thức veracity}
\label{sec:theory}

Mục này xây dựng mô hình khái niệm dùng chung cho toàn báo cáo: mọi cơ chế
trong Chương~\ref{chap:implementation} (validation per-metric, dedup,
quarantine, alert, batch aggregate) và mọi metric trong
Chương~\ref{chap:results} đều được ánh xạ về các định nghĩa ở đây. Các định
nghĩa không phụ thuộc công nghệ, nhưng mọi ví dụ cơ chế đều lấy từ pipeline
MQTT--Kafka--Spark--Parquet/PostgreSQL cụ thể của báo cáo.

Mô hình gồm hai lớp tách bạch. \textbf{Lớp baseline} là những gì đã được
hiện thực trong Chương~\ref{chap:implementation}: bộ bốn trường
định danh/thời gian, ba đại lượng độ trễ, vector veracity cùng cơ chế
per-metric, marker khử trùng và upsert idempotent ở PostgreSQL, và các bất
biến kế toán tương ứng. \textbf{Lớp extension} là những gì chỉ được định
nghĩa ở mức khái niệm, chưa triển khai và chưa có số đo: ngữ nghĩa
cửa sổ/watermark ở mức query, chính sách thích nghi FAST/CAUTIOUS và vòng
đời công bố PROVISIONAL--FINAL--CORRECTED. Lớp extension đóng vai trò khung
cho các giả thuyết cần đo (Mục~\ref{subsec:hypotheses}) và cho hướng phát
triển (Chương~\ref{chap:conclusion}); trong toàn bộ phần mô tả baseline,
báo cáo không dựa vào bất kỳ định nghĩa extension nào.

\subsection{Event, bản tin, row và vector chỉ số}
\label{subsec:event-tuple}

Trước hết, báo cáo chuẩn hoá bốn từ khoá và dùng nhất quán ở mọi chương:

\begin{itemize}
  \item \textbf{event} là đơn vị dữ liệu logic: một quan sát của một cảm
        biến tại một thời điểm, được định danh bởi \texttt{event\_id};
  \item \textbf{bản tin} (message) là biểu diễn của một event trên đường
        truyền --- một bản tin MQTT hoặc một bản ghi trong log Kafka;
  \item \textbf{row} là biểu diễn của một event sau khi được lưu trữ ---
        một dòng trong Parquet (Bronze, Silver, Quarantine) hoặc trong
        PostgreSQL;
  \item \textbf{reading} là một thành phần chỉ số của event --- một phần
        tử của vector đo --- và là đơn vị gốc của aggregate ở tầng Gold.
\end{itemize}

Sự phân biệt này quyết định ngữ nghĩa của từng cơ chế: dedup diễn ra trên
event (theo \texttt{event\_id}), quarantine diễn ra trên row, còn aggregate
và cảnh báo diễn ra trên reading. Một bản tin mang một event; một event được
lưu thành row; một event có thể sinh nhiều reading, mỗi chỉ số nó mang một
reading.

Với các định danh trên, mỗi event được mô hình hoá thành
\begin{equation}
  \begin{aligned}
    e &= \bigl(\mathit{event\_id},\ \mathit{sensor\_id},\ t^{e},\ t^{gw},\
         \mathbf{m}\bigr), \\
    \mathbf{m} &= (m_{\mathrm{temp}},\ m_{\mathrm{hum}},\ \ldots,\
         m_{\mathrm{light}}) \in \bigl(\mathbb{R} \cup \{\bot\}\bigr)^{6}.
  \end{aligned}
  \label{eq:event-tuple}
\end{equation}
theo đúng data contract khai báo trong Listing~\ref{lst:schema}: bốn trường
định danh và thời gian \textbf{bắt buộc} --- \texttt{event\_id},
\texttt{sensor\_id}, \texttt{event\_time} ($t^{e}$) và \texttt{ingest\_time}
($t^{gw}$) --- cùng vector đo $\mathbf{m}$ gồm sáu chỉ số môi trường (nhiệt
độ, độ ẩm, CO$_2$, áp suất, bụi mịn PM2.5, ánh sáng). Bốn trường này là điều
kiện cần để một row được chấp nhận vào Silver; thiếu hoặc sai một trong
chúng là lỗi định danh và được xử lý ở Mục~\ref{subsec:veracity}.

Mỗi thành phần $m_j$ nhận một giá trị số hoặc vắng mặt ($\bot$) tuỳ loại
trạm: trạm thời tiết phát nhiệt/ẩm/áp suất/ánh sáng, trạm chất lượng không
khí phát CO$_2$/PM2.5. \textbf{Chỉ số vắng mặt là dữ liệu bình thường,
không phải lỗi} (baseline đã triển khai): một row Silver hợp lệ chỉ mang tập
con chỉ số của trạm phát nó. Lỗi của từng chỉ số --- giá trị không parse
được số hoặc nằm ngoài miền vật lý --- được xử lý per-metric ở
Mục~\ref{subsec:veracity} và không làm mất row.

\subsection{Thời gian của event và ba đại lượng độ trễ}
\label{subsec:lateness}

Một event mang theo nhiều mốc thời gian; việc phân biệt chúng là nền tảng
của mọi định nghĩa độ trễ phía sau. Báo cáo dùng hệ ký hiệu thống nhất
trong Bảng~\ref{tab:time-taxonomy} ở mọi chương.

\begin{table}[H]
\centering
\caption{Phân loại các mốc thời gian của một event trong pipeline.}
\label{tab:time-taxonomy}
\begin{tabularx}{\textwidth}{c l X}
\toprule
\textbf{Ký hiệu} & \textbf{Tên} & \textbf{Điểm gắn và vai trò} \\
\midrule
$t^{e}$ & event time &
  Thời điểm hiện tượng vật lý xảy ra, do cảm biến gắn; là trục thời gian
  của mọi aggregate trong báo cáo (trường \texttt{event\_time}). \\
$t^{gw}$ & producer/gateway ingest time &
  Thời điểm producer đóng gói và phát bản tin; đại diện cho độ trễ phía
  ngoài nền tảng (trường \texttt{ingest\_time}). \\
$t^{k}$ & Kafka record timestamp &
  Timestamp gắn với bản ghi Kafka, được lưu ở Bronze dưới trường
  \texttt{kafka\_timestamp}. Baseline dùng \texttt{CreateTime} mặc định,
  do Kafka producer gắn; chỉ với cấu hình \texttt{LogAppendTime} mới
  biểu thị thời điểm broker append. Không dùng mốc này để tính độ trễ
  của báo cáo. \\
$t^{r}$ & platform receive time &
  Mốc xử lý micro-batch của query Q1, được gắn bằng
  \texttt{current\_timestamp()} và lưu thành trường provenance ở tầng
  Bronze (\texttt{received\_at\_utc}). \\
$t^{o}$ & publication time &
  Thời điểm dòng aggregate của một cửa sổ xuất hiện ở bảng Gold; tồn tại
  ở mức window, không phải mức event. \\
\bottomrule
\end{tabularx}
\end{table}

\begin{figure}[H]
\centering
\begin{tikzpicture}[
  font=\footnotesize,
  stage/.style={draw=blue!55!black, rounded corners=2pt, fill=blue!5,
    align=center, text width=2.85cm, minimum height=1.6cm, inner sep=4pt},
  link/.style={-{Latex[length=2mm]}, thick, draw=black!65},
  span/.style={{Latex[length=2mm]}-{Latex[length=2mm]}, thick, draw=blue!65!black}]
\node[stage] (event) at (0,0) {\textbf{Cảm biến}\\$t^{e}$\\\texttt{event\_time}};
\node[stage] (producer) at (3.75,0) {\textbf{Producer}\\$t^{gw}$\\\texttt{ingest\_time}};
\node[stage] (kafka) at (7.5,0) {\textbf{Kafka}\\$t^{k}$\\\texttt{kafka\_timestamp}};
\node[stage] (spark) at (11.25,0) {\textbf{Spark Q1}\\$t^{r}$\\\texttt{received\_at\_utc}};
\draw[link] (event.east) -- (producer.west);
\draw[link] (producer.east) -- (kafka.west);
\draw[link] (kafka.east) -- (spark.west);
\foreach \x in {0,3.75,11.25} {\draw[dashed, black!35] (\x,-.85) -- (\x,-3.1);}
\draw[span] (0,-1.65) -- (3.75,-1.65);
\node[align=center, text width=3.5cm] at (1.875,-2.22)
  {$t^{gw}-t^{e}$\\đóng gói và trễ phía nguồn};
\draw[span] (3.75,-1.65) -- (11.25,-1.65);
\node[align=center, text width=6.6cm] at (7.5,-2.22)
  {$t^{r}-t^{gw}$\\truyền nhận, hàng đợi và chờ micro-batch};
\draw[span] (0,-3.1) -- (11.25,-3.1);
\node[align=center, text width=12.5cm] at (5.625,-3.64)
  {\textbf{Ingestion lateness:} $L_i=\max(0,t_i^{r}-t_i^{e})$; không đồng nhất với độ trễ mạng.};
\node[draw=black!35, fill=black!3, rounded corners=2pt,
  align=left, text width=13.4cm, inner sep=5pt] at (5.625,-4.76)
  {\textbf{Lưu ý về $t^{k}$:} baseline dùng Kafka \texttt{CreateTime} mặc định.
   Timestamp do producer gắn không phải thời điểm broker append.
   $t^{r}$ là timestamp xử lý của micro-batch Q1, không phải timestamp riêng cho từng gói mạng.};
\end{tikzpicture}
\caption{Các mốc thời gian và hai thành phần của ingestion lateness trong baseline. Vị trí các khối không theo tỷ lệ thời gian; phép so sánh giả định clock đã đồng bộ.}
\label{fig:event-time-path}
\end{figure}


Hình~\ref{fig:event-time-path} cho thấy cùng một bản ghi có nhiều mốc
thời gian với ý nghĩa khác nhau. Trong baseline, Kafka giữ timestamp
\texttt{CreateTime}~\cite{kafka_docs_4_3}; trường này không đo được riêng
thời điểm ghi log ở broker. Timestamp Q1 được gắn ở mức micro-batch, nên
lateness suy ra từ Bronze bao gồm cả thời gian chờ xử lý và cần được đọc
theo độ phân giải của trigger, thay vì xem là phép đo độ trễ mạng của từng
gói.

\textbf{Giả định clock.} Các so sánh timestamp trong báo cáo giả định clock
của cảm biến/gateway và của nền tảng được đồng bộ (ví dụ NTP) hoặc độ lệch
clock (skew) bị chặn trên bởi một hằng $\varepsilon$ nhỏ so với các ngưỡng
lateness được sử dụng. Nếu giả định này bị vi phạm, mọi ước lượng lateness
đều mang sai số hệ thống cỡ $\varepsilon$ và cần được nêu rõ khi phân tích.

Trên nền các mốc thời gian đó, \textbf{lateness quan sát được} (observed
lateness, hay \textit{ingestion lateness}) tại điểm nền tảng nhận là
\begin{equation}
  L_i \;=\; \max\bigl(0,\ t^{r}_i - t^{e}_i\bigr).
  \label{eq:observed-lateness}
\end{equation}
Đây là đại lượng mức event, tính được trực tiếp từ hai cột timestamp đã lưu
trong Bronze/Silver (\texttt{received\_at\_utc} và \texttt{event\_time}),
nên thuộc trạng thái \emph{có thể suy ra từ dữ liệu đã ghi}. Thành phần
$t^{r} - t^{e}$ không đồng nhất với network delay: $L_i$ là tổng hợp của
thời gian xử lý tại cảm biến, buffering tại gateway, retry, hàng đợi mạng
và cả sai lệch clock cục bộ (theo giả định clock nêu trên). Vì vậy một
chiến lược ``giảm latency mạng'' không nhất thiết làm giảm $L_i$ nếu nguyên
nhân nằm ở gateway hoặc clock.

Cần tách bạch ngay từ đầu ba đại lượng ``độ trễ'' khác cấp độ mà báo cáo
sử dụng, không gộp chung thành một con số:

\begin{itemize}
  \item \textbf{Ingestion lateness} $L_i$ --- mức event, tính từ timestamp
        của chính event theo Ct.~\ref{eq:observed-lateness}; đặc trưng cho
        khoảng cách giữa hiện tượng và lúc nền tảng thấy nó;
  \item \textbf{Kafka consumer lag} --- mức luồng, khoảng cách giữa offset
        cuối cùng của log partition (Log End Offset) và offset mà query
        đang đọc; đặc trưng cho tích luỹ nguồn khi tốc độ vào vượt tốc độ
        xử lý;
  \item \textbf{Publication latency} --- mức window, khoảng thời gian từ
        khi cửa sổ kết thúc theo event time đến khi dòng aggregate tương
        ứng xuất hiện ở bảng Gold.
\end{itemize}
Ba đại lượng này được định nghĩa bằng công thức riêng trong
Chương~\ref{chap:results}; ở đây chỉ cần cố định tên gọi và cấp độ của
chúng.

Một event $e_j$ gọi là \textbf{đến sai thứ tự} (out-of-order) nếu tồn tại
$e_i$ với $t^{e}_i < t^{e}_j$ nhưng $t^{r}_i > t^{r}_j$. Out-of-order là hệ
quả trực tiếp của lateness không đồng đều giữa các event, và là lý do các
phép tính aggregate đặt trên trục $t^{e}$ thay vì trục thời gian đến.

Việc một event có ``đến sau watermark'' hay không là một vị từ (predicate)
khác, phụ thuộc thêm vào tiến độ đọc của query tại thời điểm event đến;
vị từ này thuộc extension chưa triển khai và được định nghĩa riêng ở
Mục~\ref{subsec:watermark}. Không suy ra nó trực tiếp từ $L_i$.

\subsection{Vector veracity và cơ chế baseline}
\label{subsec:veracity}

Đặc trưng Veracity của Big Data được báo cáo cụ thể hoá thành một vector chất
lượng 5 chiều, mỗi chiều là một câu hỏi có thể kiểm tra trên từng bản ghi
hoặc từng luồng:
\begin{equation}
  \mathcal{V}(e) = \bigl(C_{\mathrm{comp}},\ C_{\mathrm{val}},\
  C_{\mathrm{uniq}},\ C_{\mathrm{time}},\ C_{\mathrm{cons}}\bigr).
  \label{eq:veracity-vector}
\end{equation}
Bảng~\ref{tab:veracity-mapping} nêu định nghĩa từng chiều và ánh xạ tới cơ
chế baseline trong Chương~\ref{chap:implementation}; toàn bộ năm cơ chế
trong bảng thuộc trạng thái \textbf{baseline đã triển khai}.

\begin{table}[H]
\centering
\caption{Năm chiều veracity và ánh xạ tới cơ chế baseline.}
\label{tab:veracity-mapping}
\small
\begin{tabularx}{\textwidth}{l X X}
\toprule
\textbf{Chiều} & \textbf{Câu hỏi} & \textbf{Cơ chế baseline tương ứng} \\
\midrule
Completeness & Row có đủ các trường định danh/thời gian bắt buộc không? &
  Q2a/Q2b kiểm tra bốn trường bắt buộc: \texttt{event\_id},
  \texttt{sensor\_id} khác rỗng; \texttt{event\_time},
  \texttt{ingest\_time} parse được, chuẩn UTC. Vi phạm $\rightarrow$ cả
  row vào Quarantine kèm \texttt{error\_reason}; thiếu
  \texttt{event\_id} dùng key dự phòng topic--partition--offset. \\
Validity & Giá trị từng chỉ số có parse được số và nằm trong miền vật lý
  hợp lệ không? &
  Kiểm tra per-metric theo bảng cấu hình (bounds inclusive): chỉ số
  không parse được hoặc ngoài miền vật lý bị \textbf{che lọc (mask)
  thành NULL} trong Silver, reason \texttt{non\_numeric:<m>} /
  \texttt{value\_out\_of\_bounds:<m>} ghi vào \texttt{metric\_issues};
  row vẫn giữ. Chỉ số vắng (station subset) không phải lỗi; vượt ngưỡng
  cảnh báo vẫn hợp lệ. \\
Uniqueness & \texttt{event\_id} có xuất hiện nhiều lần không? &
  Khử trùng ở tầng phục vụ, không trong state streaming: Q3 ghi marker
  \texttt{processed\_events} (khoá chính \texttt{event\_id}) idempotent;
  replay tăng \texttt{duplicate\_count}; re-send cùng \texttt{event\_id}
  với vector chỉ số khác được phát hiện. Batch Gold chọn một row theo
  \texttt{ROW\_NUMBER} với thứ tự xác định cho mỗi \texttt{event\_id}. \\
Timeliness & Event đến muộn bao nhiêu? &
  Cờ \texttt{is\_late} trong Silver đánh khi $t^{r} - t^{e}$ vượt
  \texttt{LATE\_THRESHOLD\_S} (mặc định 60 giây); event muộn hợp lệ
  không bị quarantine. Vị từ post-watermark là extension
  (Mục~\ref{subsec:watermark}). \\
Consistency & Các mốc thời gian có nhất quán không? &
  Baseline bảo đảm ở mức parse: cả hai mốc \texttt{event\_time} và
  \texttt{ingest\_time} bắt buộc parse được, chuẩn hoá UTC khi ghi
  Silver; ngoài mức parse, baseline chưa có kiểm tra chéo sâu hơn giữa
  hai mốc. \\
\bottomrule
\end{tabularx}
\end{table}

Bảng trên cho thấy ba mức ngữ nghĩa phân biệt rõ trong cơ chế baseline:

\begin{itemize}
  \item \textbf{Mức row}: chỉ lỗi định danh/thời gian đưa cả row vào
        Quarantine; đây là trường hợp duy nhất một row bị loại khỏi Silver;
  \item \textbf{Mức metric}: lỗi của một chỉ số chỉ làm mất giá trị chỉ số
        đó (mask NULL kèm reason), không ảnh hưởng các chỉ số khác của cùng
        row và không đưa row vào Quarantine;
  \item \textbf{Mức event}: bản trùng vẫn hiện diện trong Silver (Silver cố
        ý cho phép lặp để không mất dữ liệu ở tầng tệp) nhưng chỉ được công
        bố một lần ở tầng phục vụ, qua marker \texttt{processed\_events} và
        dedup của batch.
\end{itemize}

\begin{figure}[H]
\centering
\begingroup
\setstretch{1}
\begin{tikzpicture}[
  font=\footnotesize,
  block/.style={draw=blue!55!black, rounded corners=2pt, fill=blue!5,
    align=center, text width=5.1cm, minimum height=1.15cm, inner sep=5pt},
  reject/.style={block, text width=3.65cm, fill=orange!9, draw=orange!65!black},
  output/.style={block, text width=4.85cm, fill=green!7, draw=green!40!black},
  link/.style={-{Latex[length=2mm]}, thick, draw=black!65}]
\node[block, text width=8.0cm] (bronze) at (0,0)
  {\textbf{Bronze}\\Mỗi bản ghi Kafka được lưu nguyên payload cùng provenance.};
\node[block] (row) at (0,-1.8)
  {\textbf{Mức row: kiểm tra định danh/thời gian}\\Bốn trường bắt buộc hợp lệ?};
\node[reject] (quarantine) at (-3.9,-3.65)
  {\textbf{Quarantine}\\Giữ cả row lỗi cùng \texttt{error\_reason}.};
\node[block] (metric) at (3.2,-3.65)
  {\textbf{Mức metric: chỉ số}\\Lỗi $\rightarrow$ NULL + \texttt{metric\_issues}.\\Chỉ số vắng không phải lỗi.};
\node[block, fill=green!7, draw=green!40!black] (silver) at (3.2,-6.0)
  {\textbf{Silver}\\Giữ row hợp lệ, các chỉ số còn dùng được và cờ \texttt{is\_late}; bản trùng vẫn được lưu.};
\node[output] (serving) at (0,-8.25)
  {\textbf{Mức event: phục vụ Q3}\\Marker \texttt{processed\_events}; upsert giá trị mới nhất và cảnh báo theo khoá.};
\node[output] (gold) at (5.7,-8.25)
  {\textbf{Mức event: batch Gold}\\Chọn một row/\texttt{event\_id} theo thứ tự xác định; tách reading; tổng hợp theo giờ.};
\draw[link] (bronze) -- (row);
\draw[link] (row.south) -- (0,-2.65) -| node[pos=.65, above]{không} (quarantine.north);
\draw[link] (row.east) -- (3.2,-1.8) -- node[right]{có} (metric.north);
\draw[link] (metric) -- (silver);
\draw[link] (silver.south) -- (3.2,-7.15) -| (serving.north);
\draw[link] (3.2,-7.15) -| (gold.north);
\node[draw=black!35, fill=black!3, rounded corners=2pt,
  text width=13.3cm, align=center, inner sep=5pt] at (1.1,-10.2)
  {\textbf{Kế toán tầng tệp sau drain, trước khử trùng:}
   $N_{\mathrm{Bronze}}=N_{\mathrm{Silver}}+N_{\mathrm{Quarantine}}$.\\
   Không cộng lại số bản trùng ở tầng phục vụ vào đẳng thức này.};
\end{tikzpicture}
\endgroup
\caption{Ba mức xử lý chất lượng dữ liệu của baseline: row quyết định nhánh lưu trữ, metric quyết định giá trị giữ lại và event quyết định tính duy nhất khi phục vụ hoặc tổng hợp.}
\label{fig:validation-levels}
\end{figure}


Hình~\ref{fig:validation-levels} làm rõ vì sao ``dữ liệu lỗi'' không đồng
nghĩa với việc bỏ cả bản tin. Một row có định danh hợp lệ vẫn vào Silver
nếu chỉ một metric lỗi; giá trị metric đó được che lọc và giữ lý do để
hậu kiểm. Cũng tại Silver, bản trùng được giữ để bảo toàn lịch sử tiếp
nhận. Batch Gold chỉ chọn một bản theo \texttt{event\_id}, ưu tiên
\texttt{received\_at\_utc} mới nhất, rồi \texttt{kafka\_partition} và
\texttt{kafka\_offset} giảm dần; sau đó chỉ các reading khác NULL mới góp
vào thống kê. Quy tắc chọn xác định này khác với việc chọn tuỳ ý một bản
trong nhóm trùng.

\subsection{Metric chất lượng rolling và trạng thái bằng chứng}
\label{subsec:rolling-metrics}

Để biến veracity vector thành đại lượng vận hành, báo cáo định nghĩa sáu
metric rolling, tính trên cửa sổ trượt (hoặc theo epoch chính sách trong
Mục~\ref{subsec:policy} khi extension đó được triển khai). Mỗi metric được
gắn trạng thái bằng chứng theo quy ước của chương:

\begin{enumerate}
  \item \textbf{p95 lateness} [\emph{có thể suy ra từ dữ liệu đã ghi}]:
        phân vị 95 của $L_i$ (Ct.~\ref{eq:observed-lateness}) trên các
        event nhận trong khoảng rolling; đặc trưng cho đuôi của phân phối
        đến muộn. Tính được từ hai cột timestamp đã lưu trong
        Bronze/Silver;
  \item \textbf{Late-event ratio} $\rho_{\mathrm{late}}$ [\emph{baseline:
        cờ đã ghi; ratio có thể suy ra}]: tỷ lệ event có $L_i$ vượt một
        ngưỡng $\ell$ chọn trước trên tổng event nhận trong khoảng.
        Baseline dùng $\ell$ = \texttt{LATE\_THRESHOLD\_S} (60 giây),
        trùng với ngưỡng của cờ \texttt{is\_late} đã ghi trong Silver.
        Ngưỡng đặt trên $L_i$ là một chỉ báo vận hành, không đồng nhất với
        vị từ post-watermark (Mục~\ref{subsec:watermark});
  \item \textbf{Duplicate ratio} $\rho_{\mathrm{dup}}$ [\emph{có thể suy ra
        từ dữ liệu đã ghi}]: trong tập row Silver được chọn, tỷ lệ bản
        thừa so với số \texttt{event\_id} phân biệt, tức
        $(N_{\mathrm{Silver}}-N_{\mathrm{distinct\ event\_id}})
        /N_{\mathrm{Silver}}$ với cùng phạm vi quan sát. Bộ đếm
        \texttt{duplicate\_count} ở Q3 còn bao gồm replay do micro-batch
        chạy lại, nên được báo riêng và không tự động bằng tỷ lệ trùng
        sinh ra ở nguồn;
  \item \textbf{Invalid/quarantine ratio} $\rho_{\mathrm{inv}}$
        [\emph{có thể suy ra từ dữ liệu đã ghi}]: tỷ lệ row bị tách sang
        Quarantine vì lỗi định danh; riêng tỷ lệ chỉ số bị che lọc (NULL
        kèm reason trong \texttt{metric\_issues}) được thống kê
        per-metric;
  \item \textbf{Watermark lag} [\emph{extension chưa triển khai}]: hiệu
        giữa đồng hồ vật lý hiện tại và watermark của query đang chạy;
        lag tăng kéo dài là dấu hiệu pipeline không theo kịp input rate.
        Ở baseline chưa khai báo watermark, backlog được quan sát bằng
        Kafka consumer lag (Mục~\ref{subsec:lateness}); consumer lag theo
        offset không phải watermark lag theo thời gian;
  \item \textbf{Policy transition count} [\emph{extension chưa triển
        khai}]: số lần chuyển trạng thái FAST $\leftrightarrow$ CAUTIOUS
        trong một khoảng quan sát; chỉ có ý nghĩa khi chính sách ở
        Mục~\ref{subsec:policy} được triển khai.
\end{enumerate}

Ba tỷ lệ (late-event, duplicate, invalid/quarantine) là ước lượng mẫu; sai số
thống kê của chúng tuân theo quy luật
lấy mẫu thông thường và cần được nêu kèm khoảng tin cậy khi dùng làm ngưỡng
quyết định --- đây là nội dung của giả thuyết H2
(Mục~\ref{subsec:hypotheses}).

\subsection{Extension event-time: cửa sổ và watermark (chưa triển khai)}
\label{subsec:watermark}

Toàn bộ mục này thuộc \textbf{lớp extension chưa triển khai}: các query
streaming của baseline không khai báo watermark (được chứng minh ở cuối
mục), nên mọi vị từ và ngữ nghĩa dưới đây mô tả hành vi của một query
giả định có khai báo watermark, dùng làm khung cho extension và cho các giả
thuyết cần đo.

\textbf{Cửa sổ event-time.} Với độ dài cửa sổ $w$ và mốc căn chỉnh $t_0$,
cửa sổ thứ $m$ là nửa khoảng
\begin{equation}
  W_m = [\,t_0 + m w,\ t_0 + (m+1) w\,)
  \label{eq:window-def}
\end{equation}
và event $e_i$ được gán vào $W_m$ khi $t^{e}_i \in W_m$ (tumbling window;
mọi window trong báo cáo thuộc loại này). Khái niệm cửa sổ trên trục
$t^{e}$ không thuộc riêng extension: batch Gold của baseline cũng aggregate
theo cửa sổ giờ trên \texttt{event\_time}. Điểm chỉ có ở extension là
watermark --- biến trạng thái quyết định ``cửa sổ khi nào được coi là đã
nhận đủ''.

\textbf{Watermark.} Với một query $Q$ đọc luồng event và ngưỡng trễ
$\theta \geq 0$, watermark được mô hình hoá cho một nguồn event-time
tại thời điểm $t$ là
\begin{equation}
  W_Q(t) \;=\; M_Q(t) - \theta,
  \qquad
  M_Q(t) = \max_{e \in E_Q(t)} t^{e}
  \label{eq:watermark-def}
\end{equation}
trong đó $E_Q(t)$ là tập event mà $Q$ đã đọc tính đến $t$. Vì $M_Q(t)$ không
giảm theo $t$, watermark có tính đơn điệu:
\begin{equation}
  t_1 \leq t_2 \;\Longrightarrow\; W_Q(t_1) \leq W_Q(t_2).
  \label{eq:watermark-mono}
\end{equation}
Đây là mô hình khái niệm; Spark cập nhật watermark theo tiến độ
micro-batch, và query nhiều nguồn còn có chính sách ghép watermark riêng.
Vì vậy không xem Ct.~\ref{eq:watermark-def} là một phép cập nhật tức thời
cho từng record trong engine.

\textbf{Lateness so với watermark.} Cho $a_i$ là thời điểm query $Q$ xử lý
record chứa $e_i$ (ở Bronze, xấp xỉ $t^{r}_i$). Event $e_i$ gọi là
\textbf{đến sau watermark} (post-watermark) của $Q$ khi tại thời điểm đến,
watermark đã vượt qua event time của nó, tức $t^{e}_i \leq W_Q(a_i)$.
Đây là vị từ khác với observed lateness $L_i$
(Ct.~\ref{eq:observed-lateness}): $L_i$ chỉ phụ thuộc hai timestamp của
chính event, còn vị từ post-watermark phụ thuộc cả tiến độ đọc $M_Q$ của
query. Hai vị từ chỉ tương quan, không tương đương: với $\theta = 10$ phút,
một event có $L_i = 12$ phút vẫn có thể đến trước watermark nếu luồng đến
chậm ($M_Q$ còn thấp), và ngược lại một event có $L_i = 6$ phút vẫn có thể
đến sau watermark nếu các event khác đã đưa $M_Q$ vượt qua event time của
nó. Do đó không được suy ``$L_i > \theta$ nên event post-watermark'' như
một kết luận tất yếu.

\textbf{Ngữ nghĩa.} Watermark là một tuyên bố tiến độ ở mức query/window: ``$Q$
xem như đã nhận đủ mọi event có $t^{e} \leq W_Q(t)$''. Nó phục vụ hai việc:
(i) giới hạn state (ví dụ state dedup chỉ cần giữ trong horizon $\theta$);
(ii) hỗ trợ \textit{finalization} --- quyết định khi nào một cửa sổ được chốt
kết quả trong chế độ streaming. Watermark là cơ chế đánh đổi
correctness--latency--cost kinh điển của xử lý dữ liệu unbounded: $\theta$ lớn
giữ được nhiều event muộn nhưng state lớn và kết quả chốt muộn; $\theta$ nhỏ
ngược lại~\cite{dataflow_model_2015}.

Hai hiểu lầm cần loại bỏ ngay từ đầu:
\begin{itemize}
  \item Watermark không phải ``thùng rác'': event đến sau watermark
        (post-watermark) \emph{không} tự động bị xoá hay tự động vào
        quarantine. Hành vi cụ thể phụ thuộc loại query và output mode của
        nó. Với query chỉ chiếu và ghi bản ghi thô, không có operator
        stateful, khai báo watermark không tự tạo một bộ lọc xoá bản ghi.
        Với operator aggregation hoặc dedup có state, bản quá muộn có
        thể bị loại theo ngữ nghĩa operator. Spark bảo đảm giữ dữ liệu
        nằm trong ngưỡng trễ khai báo, nhưng dữ liệu ngoài ngưỡng có thể
        được xử lý hoặc bị loại, không có bảo đảm một chiều ngược lại.
        Khi cửa
        sổ đã được chốt và state của nó đã được giải phóng, event
        post-watermark không thể cập nhật aggregate đó nữa --- chế độ
        \texttt{append} không phát hành thêm dòng cho cửa sổ đã chốt, và
        chế độ \texttt{update} cũng vậy vì state đã bị dọn. Chế độ
        \texttt{complete} giữ toàn bộ state aggregate và công bố toàn bộ
        bảng kết quả mỗi trigger, nên không dùng watermark để dọn state
        cũ; baseline không dùng chế độ này~\cite{spark_docs_4_2}.
  \item Watermark không phải độ trễ mạng và không phải SLA: nó là biến trạng
        thái của query, tiến theo dữ liệu thực tế đã đọc, không theo đồng hồ
        vật lý.
\end{itemize}

Trong Structured Streaming, watermark được khai báo ở mức query bằng
\texttt{withWatermark} và gắn với một cột event-time~\cite{spark_docs_4_2,
structured_streaming_sigmod_2018}. Ranh giới với cài đặt phải nêu rõ:
\textbf{không query nào trong baseline khai báo watermark}. Dữ liệu đến
muộn được xử lý qua cờ \texttt{is\_late} theo ngưỡng
\texttt{LATE\_THRESHOLD\_S} và qua tính lại batch; các query streaming
Q1--Q3 là phép chiếu, lọc và \texttt{foreachBatch} stateless nên không có
state aggregate cần chốt. Do đó $\theta$ hiện chỉ là tham số của mô hình và
của extension thích nghi (Mục~\ref{subsec:policy}), chưa phải cấu hình đang
chạy.

\subsection{Extension vận hành: chính sách thích nghi và vòng đời công bố (chưa triển khai)}
\label{subsec:policy}
\label{subsec:lifecycle}

Toàn bộ mục này cũng thuộc \textbf{lớp extension chưa triển khai}; các cơ
chế được mô tả để định nghĩa giả thuyết H4 và để chỉ ra seam nâng cấp
trong kiến trúc hiện có, không mô tả hành vi đang chạy.

\textbf{Chính sách hai trạng thái (đề xuất).} Trên nền metric rolling, báo
cáo đề xuất một chính sách hai trạng thái điều chỉnh hành vi xử lý theo
chất lượng quan sát được. Đây là thiết kế \textit{policy epoch / feedback}:
cứ sau mỗi epoch (ví dụ một cụm batch), controller đọc vector metric, so
với ngưỡng và quyết định trạng thái cho epoch kế tiếp. Chính sách thay đổi
tham số của query (ngưỡng $\theta$, horizon dedup) giữa các epoch ---
\emph{không} phải cơ chế watermark per-record native của Spark; mỗi query tại
một thời điểm vẫn chạy với một watermark cố định theo ngữ nghĩa
Mục~\ref{subsec:watermark}.

Hai trạng thái:
\begin{itemize}
  \item \textbf{FAST}: chất lượng quan sát tốt (luồng đến đúng giờ, ít trùng,
        ít lỗi). Nước đi của trạng thái này: $\theta$ nhỏ (khởi điểm từ giá
        trị đo được của p95 lateness), state dedup gọn, kết quả chốt sớm.
  \item \textbf{CAUTIOUS}: chất lượng quan sát xấu đi (p95 lateness hoặc
        $\rho_{\mathrm{late}}$ tăng). Nới $\theta$ và horizon dedup để giữ
        thêm event muộn, đánh đổi bằng state lớn hơn và publication chậm hơn.
\end{itemize}

Ký hiệu $\phi$ là metric điều khiển (ví dụ p95 lateness chuẩn hoá), $u$ là
ngưỡng lên và $d$ là ngưỡng xuống ($d < u$). Chuyển trạng thái tức thời theo
một batch duy nhất dễ gây dao động (flapping) khi metric nhiễu quanh ngưỡng,
nên báo cáo dùng \textbf{hysteresis theo số batch liên tiếp}: trạng thái chỉ
đổi khi điều kiện được giữ liên tiếp trong $n$ epoch:
\begin{equation}
  \text{FAST} \xrightarrow{\ \phi \geq u\ \text{liên tiếp}\ n\ \text{epoch}\ }
  \text{CAUTIOUS}
  \qquad
  \text{CAUTIOUS} \xrightarrow{\ \phi \leq d\ \text{liên tiếp}\ n\ \text{epoch}\ }
  \text{FAST}.
  \label{eq:hysteresis}
\end{equation}

\textbf{Hard safety bounds.} Để policy không thể tự nới lỏng hệ thống ra ngoài
miền an toàn, tồn tại các bound cứng không bị policy vượt qua, độc lập với
trạng thái: (i) $\theta \leq \theta_{\max}$ (chặn trên của watermark);
(ii) horizon dedup và thời gian giữ state có chặn trên theo dung lượng;
(iii) mọi record vẫn phải được kế toán theo bất biến I1
(Mục~\ref{subsec:invariants}) --- policy chỉ thay đổi cách phân loại đúng hạn
so với muộn, không thay đổi tổng số record phải giải trình.

Điểm can thiệp (seam) của extension trong kiến trúc Chương~\ref{chap:implementation}
là nơi controller đọc metric từ nguồn quan sát đã có (Mục quan sát vận hành
phía dưới) và khởi động lại query với tham số mới --- tức là tác động vào khâu
khởi tạo query, không thay đổi mô hình thực thi (Bronze/Silver vẫn streaming,
Gold vẫn batch). Lợi ích và chi phí của chính sách này là đối tượng của giả
thuyết H4 (Mục~\ref{subsec:hypotheses}).

\textbf{Giao thức liên tục giữa các epoch (đề xuất, chưa kiểm chứng).} Vì
watermark được tính lại từ tiến độ đọc của từng query
(Ct.~\ref{eq:watermark-def}), một query khởi động mới --- kể cả khi chỉ đổi
$\theta$ --- không tự động kế thừa watermark của query cũ: nếu state không
được khôi phục, watermark sẽ bắt đầu từ mức thấp và có nguy cơ nhảy lùi so
với mốc frontier (giá trị watermark cuối cùng) mà pipeline đã đạt được. Để
bất biến I2 (Mục~\ref{subsec:invariants}) giữ đúng xuyên các epoch, báo cáo
đề xuất giao thức tối thiểu sau:
\begin{itemize}
  \item trước khi dừng query cũ, controller chốt lại frontier watermark cuối
        cùng cùng với checkpoint/offset của query, và ưu tiên tái sử dụng cùng
        \texttt{checkpointLocation} để Spark khôi phục state (kể cả watermark)
        thay vì khởi động hoàn toàn từ đầu;
  \item frontier hiệu dụng không bao giờ bị hạ: watermark của epoch mới phải
        khởi động từ $\max(\text{frontier cũ},\ \text{watermark mới tính
        được})$;
  \item các cửa sổ biên --- window mà việc chốt kết quả phụ thuộc $\theta$ cũ
        hay mới --- phải được đối chiếu/replay bằng pass batch trước khi được
        coi là FINAL;
  \item mọi chuyển epoch được ghi audit (epoch, $\theta$ cũ/mới, frontier
        trước/sau, số cửa sổ biên) và gắn với metric policy transition count;
        hiện tượng watermark lùi được coi là vi phạm I2 cần cảnh báo.
\end{itemize}
Đây là thiết kế đề xuất, chưa được triển khai và chưa được kiểm chứng; tính
đơn điệu xuyên epoch là điều kiện phải được bảo đảm và audit, không phải hành
vi tự động của query restart.

Liên quan trực tiếp đến chính sách trên là câu hỏi về \emph{địa vị} của kết
quả được công bố sớm --- nội dung của vòng đời công bố.
\textbf{Vòng đời (đề xuất, chưa triển khai).} Khi kết quả window được công bố
trước khi
cửa sổ nhận đủ mọi event (điều bắt buộc nếu muốn có số sớm), mỗi aggregate cần
một vòng đời trạng thái rõ ràng. Gọi \textbf{logical key} của một aggregate là
\begin{equation}
  K = (W_m,\ \mathit{dimensions})
  \label{eq:logical-key}
\end{equation}
gồm cửa sổ $W_m$ (Ct.~\ref{eq:window-def}) và bộ chiều nhóm (baseline:
\texttt{sensor\_id} $\times$ chỉ số đo). Ba trạng thái:

\begin{itemize}
  \item \textbf{PROVISIONAL}: aggregate được công bố khi watermark chưa vượt
        qua $W_m$; giá trị có thể còn thay đổi.
  \item \textbf{FINAL}: watermark đã vượt qua $W_m$ cộng thêm một khoảng ân
        hạn (grace); aggregate được chốt và là nguồn cho truy vấn chuẩn.
  \item \textbf{CORRECTED}: sau FINAL, vẫn có thể xuất hiện event rất muộn
        hoặc bản chỉnh sửa (ví dụ sau khi đối chiếu với batch pass); aggregate
        được tính lại và ghi đè bản FINAL cũ, kèm tham chiếu tới phiên bản
        trước.
\end{itemize}

Hai yêu cầu đi kèm: (i) \textbf{reconciliation} --- một job batch định kỳ đối
chiếu các cửa sổ đã FINAL với dữ liệu đầy đủ hơn và phát hành CORRECTED khi
chênh lệch vượt ngưỡng; (ii) \textbf{idempotency} --- việc ghi lại cùng một
logical key $K$ phải là merge/upsert theo $K$, không phải append mù, để chạy
lại không tạo bản sao. Đây là tinh thần ``batch layer sửa lỗi cho speed
layer'' của kiến trúc Lambda~\cite{lambda_adsb_2023} áp dụng ở mức window.

\textbf{Khoảng cách với baseline.} Cài đặt hiện tại chỉ hiện thực một góc của
vòng đời này: aggregate Gold được tính bằng batch theo giờ, tính lại từ toàn
bộ dữ liệu Silver của cửa sổ và ghi bằng \emph{upsert idempotent theo khoá}
(cảm biến, giờ, chỉ số) cho các khoá mà lần chạy tính ra --- chạy lại cùng
cửa sổ trên cùng dữ liệu nguồn cho kết quả xác định, không cộng dồn;
câu lệnh upsert không xoá các
dòng nằm ngoài tập khoá của lần chạy. Các query baseline không khai báo
watermark, nên không có khái niệm ``cửa sổ đã chốt'' theo
Mục~\ref{subsec:watermark}. Phần còn thiếu so với vòng đời đề xuất là: kênh
công bố sớm PROVISIONAL từ streaming (serving realtime hiện dựa vào bảng
\texttt{sensor\_latest} theo event, không phải aggregate window), cột trạng
thái và cơ chế phát hành chênh lệch CORRECTED. Việc bổ sung đầy đủ lifecycle
là một phần của hướng phát triển (Chương~\ref{chap:conclusion}).

\subsection{Bất biến kế toán và giả thuyết kiểm chứng}
\label{subsec:invariants}
\label{subsec:hypotheses}

Một pipeline veracity-aware đúng phải giữ các bất biến sau, bất kể policy ở
trạng thái nào. Mỗi bất biến được gắn trạng thái (baseline đã triển khai
hoặc extension chưa triển khai) và, khi có, cơ chế cụ thể trong
Chương~\ref{chap:implementation}.

\begin{description}
  \item[\textbf{I1.}] \textbf{Kế toán row ở tầng tệp (baseline).} Với cùng
        tập bản ghi nguồn và sau khi Q2a/Q2b đã drain, mỗi row Bronze rơi
        vào đúng một trong hai nhánh: Silver hoặc Quarantine. Không khử
        trùng ở bước này, nên số row phải bảo toàn:
        \begin{equation}
          N_{\mathrm{Bronze}} \;=\; N_{\mathrm{Silver}}
          + N_{\mathrm{Quarantine}}.
          \label{eq:accounting}
        \end{equation}
        Cơ chế: Q2a/Q2b dùng hai predicate bổ trợ, độc quyền; lỗi
        định danh/thời gian giữ nguyên row ở Quarantine kèm lý do. Lỗi
        metric và bản trùng vẫn nằm trong Silver, nên không được cộng
        một lần nữa như row bị loại. Khi đối chiếu tầng phục vụ, dùng
        kế toán riêng cho số lượt thử ghi: lượt xử lý mới và lượt
        replay/trùng; \texttt{duplicate\_count} có thể tăng khi cùng
        micro-batch chạy lại dù số row Parquet không tăng. Đẳng thức
        trên cũng không thay thế phép đối chiếu số bản tin được phát với
        số bản ghi đã tới Kafka/Bronze.
  \item[\textbf{I2.}] \textbf{Đơn điệu watermark (extension).} Trong phạm vi một
        query \emph{có khai báo watermark}, $W_Q(t)$ không giảm theo $t$
        (Ct.~\ref{eq:watermark-mono}); đây là ngữ nghĩa của Spark Structured
        Streaming~\cite{spark_docs_4_2}. Vì baseline không khai báo watermark
        ở bất kỳ query nào, bất biến này hiện không áp dụng cho query nào
        đang chạy. Xuyên các epoch policy (restart query với $\theta$ mới),
        tính đơn điệu \emph{không} tự động được bảo đảm: query mới tính lại
        watermark từ tiến độ đọc của chính nó trừ khi state được khôi phục.
        Điều kiện để bất biến giữ đúng là giao thức liên tục ở
        Mục~\ref{subsec:policy} (không hạ frontier hiệu dụng, khôi phục/ghép
        frontier cũ); điều kiện này là đề xuất cần kiểm chứng và phải được
        audit, chưa phải hành vi đã đo của Spark.
  \item[\textbf{I3.}] \textbf{Khử trùng dựa trên marker bền vững (baseline).}
        Nhận diện bản trùng không nằm trong state streaming của Spark mà
        trong bảng marker \texttt{processed\_events} của PostgreSQL: mỗi
        \texttt{event\_id} đã được áp vào bảng phục vụ có đúng một dòng;
        replay tăng \texttt{duplicate\_count} thay vì tạo dòng mới; re-send
        cùng \texttt{event\_id} với vector chỉ số khác được phát hiện qua
        truy vấn payload-conflict. Giới hạn cần nêu rõ:
        \texttt{processed\_events} \emph{không phải bounded state} --- bảng
        này tăng theo số \texttt{event\_id} đã xử lý và không bị chặn bởi
        horizon thời gian; vai trò của nó là marker kế toán/khử trùng bền
        vững (sống sót qua restart, hậu kiểm được), không phải state dedup
        có chặn trên. Việc giới hạn state dedup theo horizon watermark là
        phương án của extension thích nghi (Mục~\ref{subsec:policy}); batch
        Gold loại trùng trên một khoảng giờ hữu hạn nên tập key của nó hữu
        hạn theo thiết kế.
  \item[\textbf{I4.}] \textbf{Khoá công bố duy nhất (baseline).} Mọi aggregate ở
        Gold có đúng một logical key $K$ (Ct.~\ref{eq:logical-key}): khoá
        (cảm biến, giờ, chỉ số). Cơ chế hiện tại:
        \texttt{gold.sensor\_hourly} ghi bằng \emph{upsert idempotent theo
        khoá} đối với các khoá xuất hiện trong lần chạy batch --- chạy lại
        cùng cửa sổ không cộng dồn. Nói chính xác: batch tính lại
        aggregate từ toàn bộ dữ liệu Silver của cửa sổ nên mọi khoá có dữ
        liệu trong cửa sổ đều được ghi đè giá trị mới, nhưng câu lệnh
        \emph{không xoá} các dòng nằm ngoài tập khoá của lần chạy; vì vậy
        đây là upsert idempotent theo batch, không phải full-replace toàn
        bảng.
  \item[\textbf{I5.}] \textbf{Hội tụ chỉnh sửa (extension).} Dưới giả định (a)
        lateness bị chặn trên bởi $\theta + $ grace và (b) dữ liệu nguồn cho
        phép replay đầy đủ, tồn tại thời điểm hữu hạn mà sau đó mọi aggregate
        CORRECTED bằng aggregate tính trên dữ liệu đầy đủ; số lần chỉnh sửa
        của mỗi window là hữu hạn. Cơ chế đề xuất: reconciliation batch
        trong Mục~\ref{subsec:lifecycle}~\cite{dataflow_model_2015}.
\end{description}

Trong năm bất biến trên, I1, I3 và I4 thuộc baseline đang chạy và được unit
test phủ trực tiếp (Chương~\ref{chap:implementation}); I2 và I5 là điều
kiện của extension, chưa có cơ chế chạy kèm.

Bên cạnh các bất biến \emph{phải} giữ, mô hình còn sinh ra các giả thuyết
\emph{có thể} kiểm chứng. Báo cáo đặt bốn giả thuyết
\textbf{chưa được kiểm chứng}
trong phạm vi báo cáo này; chúng xác định các phép đo bổ sung cần thực
hiện trên hệ thống:

\begin{itemize}
  \item[H1.] Tăng allowed lateness (watermark bound $\theta$) làm tăng tỷ lệ
        event muộn được giữ lại, nhưng đồng thời tăng kích thước state và làm
        chậm thời điểm công bố kết quả window. Cần đo bộ ba (retention, state,
        publication delay) theo $\theta$.
  \item[H2.] Các bộ đếm chất lượng (quarantine ratio, duplicate ratio, late
        ratio) bám đúng profile lỗi được inject trong sai số lấy mẫu: với
        $N$ record, sai lệch kỳ vọng giảm theo $1/\sqrt{N}$. Kiểm chứng bằng
        cách so sánh tỷ lệ quan sát với tỷ lệ inject đã biết trước.
  \item[H3.] Reconciliation (tính lại cửa sổ từ dữ liệu đầy đủ hơn) làm giảm
        sai lệch của aggregate so với giá trị đầy đủ tại các cửa sổ đủ cũ;
        mức giảm phụ thuộc vào tỷ lệ event đến muộn sau thời điểm FINAL.
  \item[H4.] Chính sách thích nghi FAST/CAUTIOUS chỉ có lợi khi sự thay đổi
        của workload/chất lượng kéo dài hơn chi phí chuyển policy (restart
        query, làm ấm lại state); nếu shift ngắn hơn chi phí chuyển, policy cố
        định với $\theta$ chọn đúng có thể tốt hơn. Kiểm chứng đòi hỏi
        ablation: (i) fixed vs adaptive, (ii) adaptive có hysteresis vs không
        có hysteresis.
\end{itemize}

Các giả thuyết H1--H4 chưa có bộ thực nghiệm đủ để xác nhận. Chúng được
giữ làm kế hoạch đánh giá extension; H2 cũng cần so sánh tỷ lệ inject với
tỷ lệ quan sát trên nhiều cỡ mẫu. Điều này không phủ nhận những số đo
baseline đã có trong Chương~\ref{chap:results}: phép chạy dài
ngày 02/10/2026 là số đo runtime; fixture chức năng có
evidence riêng. Các đối chứng chưa chạy được giữ dưới dạng kế hoạch
kiểm chứng, chưa phải kết quả thực nghiệm.

\section{Các công nghệ chính}

\subsection{Apache Kafka --- distributed ingestion layer}

Apache Kafka là nền tảng phân tán lưu trữ và truyền tải event theo mô hình
append-only log~\cite{kafka_docs_4_3}. Trong báo cáo, Kafka đóng vai trò
buffer giữa lớp sinh dữ liệu (sensor simulator) và lớp xử lý (Spark), cho
phép hai bên vận hành độc lập về tốc độ.

Hai khái niệm then chốt trong Kafka là \textbf{topic} và \textbf{partition}.
Một topic được phân chia thành nhiều partition; mỗi partition là một log
có thứ tự, là đơn vị nhỏ nhất của song song~\cite{confluent_partition_count}.
Số partition quyết định mức song song tối đa của consumer trong cùng một
consumer group, đồng thời ảnh hưởng trực tiếp đến throughput producer.

Từ Kafka 4.0, chế độ ZooKeeper đã bị loại bỏ; dòng Kafka 4.x vận hành
bằng KRaft~\cite{kafka_docs_4_3}. Vì vậy Docker
image chính thức \texttt{apache/kafka:4.3.1} có thể khởi động một broker
chỉ với biến môi trường mà không cần triển khai ZooKeeper riêng.

Kafka hỗ trợ các lựa chọn về ngữ nghĩa giao hàng và xử lý, gồm
\textit{at-most-once}, \textit{at-least-once} và các cơ chế
\textit{exactly-once} trong phạm vi giao dịch Kafka. Producer idempotent
tránh bản trùng do retry trong phiên producer; transaction còn cần phối
hợp dữ liệu đầu ra và consumer offset khi xây dựng xử lý exactly-once
giữa các topic Kafka~\cite{kafka_docs_4_3}. Các cấu hình này không tự tạo
giao dịch nguyên tử giữa MQTT, Parquet và PostgreSQL.
Trong kiến trúc của báo cáo, Kafka producer không phải thiết bị cảm biến
(chúng chỉ nói MQTT) mà là Bridge: Bridge dùng \texttt{acks=all} và chỉ xác
nhận lại với broker MQTT sau khi Kafka giao thành công. Cơ chế này cho
phép phát lại trong giới hạn phiên bền, dung lượng hàng đợi và dữ liệu
được broker giữ; phần trùng lặp do phát lại được khử ở hạ lưu theo định
danh sự kiện. Với một broker và replication factor bằng 1, \texttt{acks=all}
không tạo thêm bản sao hay khả năng chịu hỏng đĩa.

\begin{figure}[H]
\centering
\begingroup
\setstretch{1}
\begin{tikzpicture}[
  font=\footnotesize,
  block/.style={draw=blue!55!black, rounded corners=2pt, fill=blue!5,
    align=center, text width=3.15cm, minimum height=.85cm, inner sep=4pt},
  task/.style={block, fill=orange!7, draw=orange!55!black},
  slot/.style={block, fill=green!7, draw=green!40!black},
  link/.style={-{Latex[length=2mm]}, thick, draw=black!65}]
\node[block, minimum height=2.9cm] (keys) at (0,-1.15)
  {\textbf{Bản tin MQTT}\\key: \texttt{sensor\_id}\\A, B, C, A, \ldots\\\emph{Tên key minh hoạ}};
\node[block, minimum height=2.0cm] (producer) at (4.6,-1.15)
  {\textbf{Bridge / Producer}\\Gán partition theo key và cấu hình partitioner.};
\node[block] (p0) at (9.2,0) {Partition P0\\log có thứ tự riêng};
\node[block] (p1) at (9.2,-1.15) {Partition P1\\log có thứ tự riêng};
\node[block] (p2) at (9.2,-2.3) {Partition P2\\log có thứ tự riêng};
\node[font=\footnotesize\bfseries] at (9.2,.9) {Topic \texttt{sensor\_raw}: 3 partition};
\draw[link] (keys) -- (producer);
\draw[link] (producer.east) -- (7,-1.15) |- (p0.west);
\draw[link] (producer.east) -- (p1.west);
\draw[link] (7,-1.15) |- (p2.west);
\node[align=center, text width=12.8cm] at (4.6,-3.6)
  {\textbf{Đọc nguồn Kafka của Q1 trong một micro-batch}\\Tác vụ bám partition; số partition không chứng minh các key được chia đều.};
\node[task] (t0) at (0,-4.8) {Task đọc P0};
\node[task] (t1) at (4.6,-4.8) {Task đọc P1};
\node[task] (t2) at (9.2,-4.8) {Task đọc P2};
\node[slot] (s0) at (2.3,-6.6) {Slot thực thi 1\\một task đang chạy};
\node[slot] (s1) at (6.9,-6.6) {Slot thực thi 2\\một task đang chạy};
\draw[link] (t0.south) -- (s0.north);
\draw[link] (t1.south) -- (s1.north);
\draw[link, dashed] (t2.south) -- node[pos=.5, fill=white, inner sep=2pt, anchor=west, xshift=5pt, align=left]{chờ slot\\trống} (s1.north);
\node[draw=black!35, fill=black!3, rounded corners=2pt,
  text width=13.1cm, align=center, inner sep=5pt] at (4.6,-8.15)
  {\texttt{local[2]} cung cấp tối đa hai slot task cho SparkSession.
   Bốn query streaming dùng chung tài nguyên này. Tăng partition không tự tăng CPU, RAM hay tốc độ ghi sink.};
\end{tikzpicture}
\endgroup
\caption{Phân hoạch Kafka và giới hạn đồng thời của Spark trong cấu hình một máy. Gán key và thứ tự các task là sơ đồ minh hoạ, không phải phân bố partition hoặc lịch thực thi đã đo.}
\label{fig:partition-concurrency}
\end{figure}


Hình~\ref{fig:partition-concurrency} tách hai giới hạn thường bị nhầm:
partition tạo cơ hội đọc song song, còn số slot Spark quyết định số task
đồng thời có thể thực thi. Một key ổn định thường đi vào cùng partition
theo partitioner khi số partition không đổi, nhưng không bảo đảm các key
được chia đều. Spark Structured Streaming đọc Kafka theo tác vụ của
micro-batch; không áp dụng máy móc quy tắc consumer group để suy ra số
query có thể chạy. Trong baseline, topic có ba partition và bốn query
dùng chung SparkSession \texttt{local[2]}, nên phép đo mở rộng cần theo
dõi thêm CPU, backlog, thời lượng trigger và sink.

\subsection{Apache Spark Structured Streaming --- stream processor}

Spark Structured Streaming là một engine xử lý stream dựa trên Spark SQL
và bộ API DataFrame cùng Dataset~\cite{structured_streaming_sigmod_2018}. Khác với
các API streaming cũ (DStream) hoạt động trên RDD, Structured Streaming
mô hình hoá dữ liệu unbounded như một bảng liên tục được mở rộng theo
event-time~\cite{spark_docs_4_2}.

Hệ thống mặc định vận hành theo cơ chế \textbf{micro-batch}: dữ liệu
đến được đệm trong một khoảng trigger interval, sau đó được xử lý như
một mini-batch~\cite{spark_docs_4_2}. Nhờ vậy, cùng một mã nguồn có
thể chạy cho cả batch lẫn streaming, đơn giản hoá đáng kể quá trình phát
triển.

Engine cung cấp sẵn các primitive event-time: \textbf{event-time window}
(phép tính theo \texttt{event\_time} thay vì thời điểm nhận), \textbf{watermark}
(tiến độ ở mức query giới hạn trễ cho phép của dữ liệu muộn, hỗ trợ giới hạn
state và finalization) và \textbf{streaming deduplication} (loại trùng theo
khoá, có thể giới hạn bởi watermark)~\cite{spark_docs_4_2}; định nghĩa hình
thức của các primitive này nằm ở Mục~\ref{subsec:watermark}. Cách baseline
sử dụng cần nêu rõ để tránh suy diễn quá mức: bốn query streaming (Q1--Q3)
là phép chiếu, lọc và \texttt{foreachBatch} \emph{stateless} --- không thực
hiện aggregate, không khai báo watermark, không dùng dedup state của engine;
duy nhất batch Gold thực hiện aggregate theo event-time, tính theo giờ trên
\texttt{event\_time} sau khi loại trùng theo \texttt{event\_id}. Các
primitive nêu trên là nền tảng mà extension event-time
(Mục~\ref{subsec:watermark}, Mục~\ref{subsec:policy}) dự kiến dùng khi được
triển khai.

Về mặt fault tolerance, Structured Streaming ghi lại trạng thái của
mỗi query vào một \textit{checkpoint location}. Khi driver restart,
query được khôi phục từ checkpoint và tiếp tục từ vị trí batch cuối
cùng~\cite{spark_docs_4_2}. File sink Parquet của Structured Streaming
có metadata commit theo batch và bảo đảm \textbf{exactly-once} trong phạm
vi query với nguồn có thể replay và checkpoint được giữ. Đây không phải
giao dịch ACID đa bảng của bản thân định dạng Parquet, và cũng không phải
giao dịch nguyên tử giữa các query hoặc sink~\cite{spark_docs_4_2,
structured_streaming_sigmod_2018}. Với sink ngoài qua
\texttt{foreachBatch}, bảo đảm này không tự có --- nếu micro-batch được chạy
lại, sink bên ngoài có thể nhận bản trùng --- và phụ thuộc tính idempotent
của từng sink. Baseline của báo cáo chọn hướng idempotency ở tầng phục vụ:
sink PostgreSQL đạt được bằng staging + upsert theo khoá chính (như đã nêu
ở mục lưu trữ phía dưới), không thêm tầng table format, và báo cáo
\emph{không tuyên bố} exactly-once xuyên toàn pipeline; mục tiêu đã hiện
thực là replay an toàn và không cộng dồn ở các bảng phục vụ.

\subsection{Lưu trữ: các lớp Parquet cục bộ và PostgreSQL}

Apache Parquet là định dạng lưu trữ cột (columnar) nén, tối ưu cho truy
vấn phân tích và là định dạng đầu ra mặc định của Spark Structured
Streaming~\cite{spark_docs_4_2}. Trong bản triển khai một máy của báo cáo,
ba lớp dữ liệu Bronze/Silver/Quarantine được ghi dưới dạng Parquet trên một
thư mục bind mount cục bộ (gắn vào container tại \texttt{/data/bigdata}),
thay vì đặt trên một hệ thống tệp phân tán như HDFS~\cite{hdfs_docs_3_5}
hay một open table format như Delta Lake~\cite{delta_lake_docs}. Lựa chọn
này giữ ngân sách tài nguyên thấp cho bản demo mà vẫn bảo toàn \emph{hình
dạng} kiến trúc medallion: dữ liệu thô (Bronze), dữ liệu đã làm sạch
(Silver), dữ liệu lỗi bị cách ly (Quarantine) và dữ liệu tổng hợp (Gold)
vẫn tách biệt theo đúng vai trò.

Điểm đánh đổi so với Delta Lake/HDFS là: Parquet cục bộ không cung cấp
transaction log, ACID đa người ghi, time travel hay replication ở tầng tệp.
Vì vậy tính idempotency (chạy lại không tạo bản sao, không cộng dồn kết
quả) được đảm bảo ở \emph{tầng phục vụ} bằng các ràng buộc khoá chính và
câu lệnh upsert của PostgreSQL, cùng batch Gold tính lại từ Silver rồi ghi
bằng upsert theo khoá (chi tiết ngay dưới). Việc nâng cấp sang Delta
Lake trên HDFS hoặc object storage --- để có ACID/time travel ở tầng tệp và
replication chịu lỗi --- được ghi nhận là hướng phát triển
(Chương~\ref{chap:conclusion}) chứ không thuộc phạm vi bản một máy.

Một vấn đề thường gặp với streaming là \textit{small-file problem}: mỗi
micro-batch có thể tạo một tệp Parquet rất nhỏ, làm giảm hiệu năng truy
vấn về sau. Trong bản triển khai này, việc đọc luồng tệp (file stream) từ
Bronze/Silver được giới hạn số tệp mỗi trigger để kiểm soát áp lực, và tác
vụ gom tệp/retention có thể được đặt lịch bằng cron khi dữ liệu đủ lớn.

Nếu Parquet gánh phần lưu trữ phân tích thì phần phục vụ truy vấn dùng một
hệ quản trị quan hệ. PostgreSQL là hệ quản trị cơ sở dữ liệu quan hệ mã
nguồn mở, được dùng
trong báo cáo làm nơi lưu dữ liệu phục vụ truy vấn real-time (giá trị mới
nhất theo cảm biến, bản tin cảnh báo) và aggregate Gold theo giờ. Một
instance PostgreSQL duy nhất với database \texttt{app} phục vụ cả hai vai trò
này; các bảng dữ liệu và bảng aggregate nằm trong các schema tách biệt.

Vai trò then chốt của PostgreSQL trong kiến trúc này là gánh phần
\emph{idempotency} mà Parquet cục bộ không cung cấp: bảng giá trị mới nhất
dùng khoá chính theo cảm biến và chỉ ghi đè khi event mới hơn (so sánh
hàng giá trị có thứ tự xác định); bảng cảnh báo dùng khoá chính kép
(event, quy tắc) với xung đột thì bỏ qua; bảng marker
\texttt{processed\_events} ghi nhận mọi event đã áp vào bảng phục vụ
(replay tăng bộ đếm trùng thay vì tạo dòng mới); bảng Gold dùng
\emph{upsert idempotent theo khoá} (cảm biến, giờ, chỉ số) đối với các khoá
mà batch tính ra trong cửa sổ, nên chạy lại cùng cửa sổ không cộng dồn ---
và, như đã nêu, câu lệnh không xoá các dòng nằm ngoài tập khoá của lần chạy.
Các câu lệnh upsert này là nguồn sự thật duy nhất (single source of truth)
cho ngữ nghĩa chống trùng và được kiểm thử độc lập bằng SQLite trong bộ
unit test (Chương~\ref{chap:implementation}). Transaction giữa việc ghi
Parquet và việc ghi PostgreSQL không nguyên tử xuyên hai sink; bản ghi tạm
(staging) còn sót lại do crash giữa hai bước sẽ được dọn ở đầu lần ghi kế
tiếp, và việc đối chiếu dựa vào chính các upsert idempotent nêu trên.

\subsection{MQTT và Bridge --- lớp tiếp nhận thiết bị}

MQTT là giao thức nhắn tin nhẹ theo mô hình publish/subscribe, được thiết
kế cho thiết bị IoT có tài nguyên hạn chế và mạng không ổn định, với ba
mức chất lượng dịch vụ (QoS 0/1/2)~\cite{mqtt_v311_spec}. Trong báo cáo,
broker Mosquitto tiếp nhận bản tin QoS 1 từ các simulator cảm biến. Việc
dùng MQTT ở biên thiết bị --- thay vì cho Spark đọc thẳng từ một giao thức
khác --- tách bạch \emph{giao thức thiết bị} khỏi \emph{giao thức xử lý}
(Kafka): engine xử lý không cần biết MQTT tồn tại, nên việc thay broker,
thêm bảo mật hay đổi nguồn phát đều không lan sang tầng xử lý.

Giữa MQTT và Kafka là một \textbf{Bridge} Python đóng vai trò lớp dịch giao
thức độc lập. Bridge đăng ký MQTT với phiên bền (persistent session) và cơ
chế xác nhận thủ công (manual ACK): nó chỉ gửi ACK cho broker \emph{sau khi}
Kafka báo đã ghi bản tin thành công (\texttt{acks=all}). Nếu Bridge ngừng
giữa lúc produce và ACK, broker sẽ phát lại bản tin chưa được xác nhận khi
kết nối lại; chừng nào broker còn giữ được bản tin trong phiên bền, đoạn
tiếp nhận đạt ngữ nghĩa at-least-once và phần trùng lặp do phát lại được
khử ở hạ lưu theo định danh sự kiện. Khi hàng đợi trong bộ nhớ của Bridge
đầy, bản tin bị bỏ qua mà \emph{không} ACK để broker giữ lại và phát lại;
tuy nhiên, hàng đợi của phiên bền phía broker có giới hạn cấu hình
(\texttt{max\_queued\_messages} của Mosquitto, đặt 5000 trong bản demo), nên
khi tắc nghẽn kéo dài broker cuối cùng buộc phải bỏ bản tin mới đến --- cửa
sổ mất dữ liệu này được ghi nhận tường minh trong mã nguồn Bridge, vì thế
bảo đảm at-least-once chỉ có giá trị trong giới hạn hàng đợi phía broker
chứ không tuyệt đối. Lớp Bridge cũng là ranh giới mở rộng ngang ở tầng tiếp
nhận: về nguyên tắc, khi tải tăng có thể nhân bản nhiều instance Bridge,
mỗi instance cần một định danh phiên MQTT và một định danh producer Kafka
riêng để các luồng ghi không can thiệp lẫn nhau. Bản triển khai hiện tại
vẫn dùng một định danh cố định cho cả hai phía và chưa có cấu hình định
danh theo instance, nên nhân bản Bridge mới là hướng mở rộng được ghi nhận
ở Chương~\ref{chap:conclusion} chứ chưa phải capability đã có sẵn. Ngoài
định danh, các instance phải có cơ chế chia subscription phù hợp (ví dụ
chia topic hoặc shared subscription); các subscriber độc lập cùng đăng
ký toàn bộ topic sẽ nhận nhiều bản sao, không tự chia tải.

\begin{figure}[H]
\centering
\begin{tikzpicture}[
  font=\footnotesize,
  actor/.style={draw=blue!55!black, rounded corners=2pt, fill=blue!5,
    align=center, text width=2.65cm, minimum height=.8cm, inner sep=4pt},
  message/.style={-{Latex[length=2mm]}, thick, draw=blue!65!black},
  ack/.style={-{Latex[length=2mm]}, thick, draw=green!45!black}]
\node[actor] at (0,0) {Simulator\\publisher MQTT};
\node[actor] at (3.75,0) {Mosquitto\\broker MQTT};
\node[actor] at (7.5,0) {Bridge\\subscriber + producer};
\node[actor] at (11.25,0) {Kafka\\broker};
\foreach \x in {0,3.75,7.5,11.25} {\draw[dashed, black!40] (\x,-.5) -- (\x,-6.1);}
\draw[message] (0,-1.1) -- node[above]{1. PUBLISH QoS 1} (3.75,-1.1);
\draw[ack] (3.75,-1.95) -- node[above]{2. PUBACK cho publisher} (0,-1.95);
\draw[message] (3.75,-2.8) -- node[above]{3. PUBLISH QoS 1} (7.5,-2.8);
\draw[message] (7.5,-3.65) -- node[above]{4. produce, \texttt{acks=all}} (11.25,-3.65);
\draw[ack] (11.25,-4.5) -- node[above]{5. delivery callback: OK} (7.5,-4.5);
\draw[ack] (7.5,-5.35) -- node[above]{6. PUBACK thủ công} (3.75,-5.35);
\node[draw=orange!60!black, fill=orange!7, rounded corners=2pt,
  text width=13.3cm, align=left, inner sep=5pt] at (5.625,-6.95)
  {\textbf{Khoảng có thể phát lại:} Bridge ngừng sau khi Kafka nhận bản tin nhưng trước bước 6.
   Mosquitto giữ bản chưa ACK trong phiên bền và có thể phát lại khi Bridge kết nối lại;
   định danh \texttt{event\_id} giúp hạ lưu nhận diện bản trùng.};
\end{tikzpicture}
\caption{Hai vòng xác nhận MQTT QoS 1 độc lập và điểm ACK thủ công của Bridge. Đây là một trình tự minh hoạ; PUBACK cho publisher không xác nhận rằng Kafka hoặc PostgreSQL đã ghi thành công.}
\label{fig:mqtt-ack-boundary}
\end{figure}


Hình~\ref{fig:mqtt-ack-boundary} phân biệt PUBACK mà simulator nhận từ
Mosquitto với PUBACK Bridge gửi về Mosquitto. Hai phiên QoS 1 này độc lập:
broker có thể xác nhận publisher trước khi Bridge hoặc Kafka hoàn tất
xử lý~\cite{mqtt_v311_spec}. Thứ tự bắt buộc trong Bridge là delivery
Kafka thành công trước ACK thủ công, còn các bước xác nhận publisher và
chuyển bản tin cho subscriber có thể xen kẽ. Vì thế số bản tin simulator
được ACK không tự chứng minh số bản ghi đã tới Bronze; cần đối chiếu log
phát, log delivery và row lưu trữ theo cùng phạm vi.

\subsection{Điều phối batch và quan sát vận hành bằng script}

Tác vụ tổng hợp Gold theo giờ có điểm bắt đầu/kết thúc rõ ràng và chạy lặp
lại theo lịch. Trong bản một máy, việc này được giao cho script wrapper
\texttt{run-batch-hourly} được cron của host khởi theo chu kỳ. Script trước
hết dùng một khoá loại trừ dựa trên thư mục kèm ghi PID: từ chối lần chạy
mới nếu batch trước còn chạy và tự dọn khoá cũ nếu process đó đã chết. Sau
đó script xác định cửa sổ giờ UTC vừa kết thúc, chờ một khoảng nới lỏng cấu
hình được (\texttt{BATCH\_GRACE\_S}) để các event sát ranh giới kịp nằm
trong Silver, rồi gọi đúng chương trình batch cho cửa sổ đó --- song song
với chế độ live đang chạy, không dừng nguồn phát và không đợi pipeline
drain. Ngược lại, đường batch một lượt (\texttt{run-mode.sh batch}) dành
cho tái xử lý thủ công mới dừng simulator và Bridge, chờ Spark tiêu hết
backlog Kafka, rồi mới dừng stream job và chạy batch. Một nền tảng điều phối
tác vụ đầy đủ (Airflow hoặc tương đương) chỉ trở nên cần thiết khi số lượng
DAG, phụ thuộc giữa tác vụ và nhu cầu quan sát lịch sử chạy tăng lên; lựa
chọn đó được ghi nhận là hướng phát triển ở Chương~\ref{chap:conclusion}.

Cùng triết lý ``dùng nguồn số liệu sẵn có thay vì dựng thêm hệ thống'' là
quan sát vận hành. Trong bản triển khai một máy, hệ thống không dựng một cụm
giám sát riêng kiểu Prometheus, Grafana hay JMX Exporter (để giữ ngân sách
bộ nhớ), mà thu các luồng bằng chứng độc lập với nhau: (i) tiến độ truy vấn
của Spark Structured Streaming qua query progress API --- ứng dụng ghi mỗi
trigger một dòng JSONL kèm \texttt{numInputRows}, thời lượng trigger và
offset đầu/cuối của nguồn Kafka~\cite{spark_docs_4_2}; (ii) các bộ đếm phía
PostgreSQL/API: số event đã áp vào bảng phục vụ, số bản trùng lặp, số cảnh
báo và thời điểm xử lý gần nhất; (iii) thống kê tài nguyên container qua
\texttt{docker stats} cùng RAM và dung lượng trống của host; và (iv) trạng
thái backlog: end offset của Kafka, số tệp Parquet ở từng tầng và số dòng
các bảng PostgreSQL. Một script thu thập lấy mẫu các nguồn (iii)--(iv) theo
chu kỳ và ghi mỗi mẫu thành một dòng \texttt{resource-report.jsonl} trên
bind mount. Các luồng này là bằng chứng cho người vận hành và nguyên liệu
phân tích benchmark (Chương~\ref{chap:results}); chúng không phải kênh dữ
liệu của dashboard --- dashboard đọc trực tiếp các endpoint ứng dụng của
backend.

Cách tiếp cận này đánh đổi khả năng cảnh báo và trực quan hoá tập trung của
một bộ Prometheus cùng Grafana lấy sự gọn nhẹ phù hợp bản demo; việc chuyển
sang một cụm
quan sát đầy đủ (Prometheus pull model, JMX Exporter cho JVM của Kafka và
Spark,
Grafana alerting) được ghi nhận là hướng phát triển
(Chương~\ref{chap:conclusion}) và không làm thay đổi các metric cần thu thập,
vì các đại lượng này đều trích được từ query progress API, thống kê
container và các bộ đếm PostgreSQL. Nguồn
metric này cũng là đầu vào mà chính sách thích nghi đề xuất ở
Mục~\ref{subsec:policy} đọc vào mỗi epoch.

\section{Các nghiên cứu liên quan}
\label{sec:related-work}

Ba khảo sát tổng quan cung cấp bối cảnh cho toàn chương:
\cite{iot_big_data_survey_2022} vẽ bức tranh thách thức
ingestion--storage--processing của IoT Big Data (đã dùng ở Mục~2.1),
\cite{stream_processing_evolution_2024} tổng hợp sự tiến hoá của stream
processing theo bốn trục state management, fault tolerance, elasticity và
load management, còn \cite{iot_stream_latency_2021} phân loại các loại
latency trong truy vấn IoT stream --- cơ sở để
Chương~\ref{chap:results} phân biệt các loại trễ thay vì gộp chung một con
số. Phần còn lại của mục này so sánh các công trình trực tiếp liên quan đến
mô hình và phương pháp của báo cáo theo bốn nhóm: nền tảng event-time,
thích nghi và chỉnh sửa, thực hành benchmark, và giám sát chất lượng
stream-first.

\subsection{Nền tảng xử lý event-time}

\textbf{The Dataflow Model}~\cite{dataflow_model_2015} là công trình nền tảng
của xử lý event-time hiện đại: định nghĩa thống nhất what/where/when/how cho
pipeline trên dữ liệu unbounded, out-of-order; hệ thống hoá trigger,
accumulation mode và watermark như công cụ đánh đổi correctness--latency--cost.
\textbf{MillWheel}~\cite{millwheel_2013} đưa ra trước đó khái niệm
\textit{low watermark} cho exactly-once stream processing quy mô Internet và
cơ chế persistent state chịu lỗi. Hai công trình này cung cấp \emph{ngôn ngữ
và ngữ nghĩa} mà báo cáo kế thừa ở Mục~\ref{sec:theory}; báo cáo không nhận
event-time hay watermark là đóng góp mới của mình.

Trên nền hai công trình đó, bài báo gốc về Structured
Streaming~\cite{structured_streaming_sigmod_2018}
giải thích tại sao Spark chọn mô hình relational API cho streaming và
chứng minh tính đúng đắn của cơ chế micro-batch; đây là nguồn định nghĩa
query-level watermark/dedup mà mô hình của báo cáo sử dụng. Đóng góp của
báo cáo chỉ là cách áp dụng cụ thể các primitive này cho contract vector đa
chỉ số, không phải bản thân cơ chế.

\subsection{Thích nghi runtime và chỉnh sửa batch}

Strider~\cite{strider_iswc_2017} là engine xử lý RDF stream phân tán có khả
năng thích nghi: hệ thống theo dõi tải và có thể chuyển đổi chiến lược thực
thi (hybrid reasoning) tại runtime. Báo cáo tham khảo Strider như bằng chứng
rằng \textit{runtime adaptation trong stream processing là hướng nghiên cứu có
thật và không miễn phí}: mọi cơ chế tự chuyển đổi đều cần điều kiện kích hoạt,
chi phí chuyển và đánh giá lợi ích. Đề xuất FAST/CAUTIOUS
(Mục~\ref{subsec:policy}) kế thừa tinh thần đó nhưng ở mức đơn giản hơn nhiều
--- chuyển tham số giữa các epoch kèm hysteresis --- và chưa được triển khai
hay đo lường.

Ở góc độ khác của cùng vấn đề chất lượng, nhóm công trình thứ hai xử lý sai
lệch bằng \emph{tính lại} thay vì thích nghi khi chạy. Nghiên cứu đánh giá
hiệu năng kiến trúc Lambda trên dữ liệu giám sát hàng
không ADS-B~\cite{lambda_adsb_2023} minh hoạ vai trò của batch layer trong
việc đối chiếu và chỉnh sửa kết quả của speed layer, chấp nhận kết quả sớm
ở trạng thái chưa đầy đủ (eventual consistency về chất lượng). Đây là cơ sở
thực tiễn cho vòng đời PROVISIONAL $\rightarrow$ FINAL $\rightarrow$ CORRECTED
(Mục~\ref{subsec:lifecycle}). Điểm khác biệt: công trình ADS-B vận hành kiến
trúc Lambda đầy đủ cho một miền ứng dụng cụ thể, còn báo cáo này đề xuất một
framework veracity-aware với metric và policy gắn vào pipeline lakehouse một
luồng (không nhân đôi batch layer đầy đủ), và mới dừng ở mức khái niệm.

\subsection{Thực hành benchmark cho hệ thống event-streaming}

Một tổng hợp hệ thống gần đây về các hệ thống event-streaming dựa trên Apache
Kafka~\cite{kafka_patterns_benchmark_2025} phân loại chín mẫu thiết kế lặp lại
(log compaction, CQRS bus, pipeline exactly-once, change-data capture,
stream--table join, saga orchestrator, topic đa người thuê, lưu trữ phân tầng và
event sourcing replay) và khảo sát cách bốn mươi hai nghiên cứu đánh giá chúng
thực nghiệm. Phát hiện then chốt liên quan trực tiếp đến phương pháp của báo
cáo này là sự \emph{thiếu nhất quán trong công bố cấu hình}: ngay cả các nghiên
cứu dùng bộ chuẩn TPCx-Kafka cũng thường bỏ sót số partition, replication factor
hay kích thước thông điệp, khiến kết quả không thể so sánh khách quan giữa các
công trình. Từ đó, công trình khuyến nghị một danh mục kiểm tra công bố tối
thiểu (partition, replication factor, kích thước thông điệp, dạng tải) như điều
kiện tiên quyết cho đánh giá tái lập được. Báo cáo này áp dụng nguyên tắc đó ở
Chương~\ref{chap:results}: mọi bảng số liệu đều ghi tường minh cấu hình đo để
tránh nhầm lẫn do so sánh lệch điều kiện. Ngoài ra, các thí nghiệm trong công
trình cho thấy thông lượng tăng gần tuyến tính theo số partition và phía tiêu
thụ tăng tới $\min(\text{số consumer}, \text{số partition})$ rồi bão hoà --- dạng
đường cong mà báo cáo dùng làm kỳ vọng tham chiếu cho kịch bản suy rộng
partition (RQ-B), còn số liệu cụ thể vẫn phải là số đo thật trên máy mục tiêu.

Hai benchmark tham chiếu cụ thể được dùng trực tiếp cho thiết kế thực
nghiệm của Chương~\ref{chap:results}. \cite{shufflebench_2024} là benchmark
2024 đánh giá chi phí shuffle và hành vi throughput/latency của các framework
stream phân tán (Spark, Flink, Kafka Streams) trên testbed chuẩn hoá; báo
cáo vay phương pháp chọn biến độc lập, ramp input rate và dùng phân vị
p95/p99 từ công trình này. \cite{fault_recovery_stream_2024} dùng failure
injection để đo recovery time và event latency của Spark Structured
Streaming, Flink và Kafka Streams khi worker bị kill; kết quả cho thấy Spark
phục hồi ổn định nhờ checkpoint nhưng có đánh đổi latency trong giai đoạn
restart. Báo cáo sử dụng cơ sở này để thiết kế kịch bản fault injection.

\subsection{Giám sát chất lượng stream-first}

Trong khi các công cụ chất lượng dữ liệu truyền thống vận hành theo mô hình
tĩnh (đánh giá một lần trên dữ liệu đã lưu) hoặc tăng dần (coi streaming như
mở rộng của batch trên một tập dữ liệu rồi sẽ ``hoàn chỉnh''), một luồng dữ
liệu vô hạn không bao giờ đạt trạng thái hoàn chỉnh và các vấn đề chất lượng
có thể thoáng qua nhưng tác động lớn~\cite{stream_daq_2024}. Công trình
Stream DaQ đề xuất một mô hình giám sát chất lượng thiết kế riêng cho luồng:
đánh giá liên tục trên các cửa sổ thời gian cấu hình được, với năm nhóm phép
kiểm tra --- theo từng bản ghi (tuple-at-a-time), theo ngữ cảnh cửa sổ, theo dữ
liệu tham chiếu, thích nghi động theo ngữ cảnh, và theo khoá (keyed, chia luồng
theo một thuộc tính như định danh cảm biến) --- cùng khả năng tổ hợp chúng
thành các ràng buộc phức tạp. Một khái niệm trung tâm là \emph{quality
meta-stream}: biến kết quả giám sát thành một luồng có cấu trúc
$\langle t_{\text{start}}, t_{\text{end}}, \text{measurement}, \text{assessment}\rangle$
chảy song song, dùng làm tín hiệu nhận thức chất lượng cho toàn pipeline.

Công trình này liên quan trực tiếp đến hai lựa chọn thiết kế của báo cáo. Thứ
nhất, phép kiểm tra theo từng bản ghi (tuple-at-a-time) chính là cách lớp Silver
đánh giá độc lập mỗi event theo nhiều ràng buộc, và việc mở rộng sang sáu chỉ
số môi trường với miền giá trị/đơn vị/hướng ngưỡng khác nhau buộc các phép kiểm
tra phải tổng quát theo cấu hình chứ không viết cứng cho một đại lượng --- nếu
không, sai sót sẽ lộ ngay ở các chỉ số có đặc tính khác nhiệt độ. Thứ hai, lý
do cần nhiều loại chỉ số và reason lỗi có cấu trúc (ví dụ
\texttt{value\_out\_of\_bounds:co2}) là để tránh nhầm lẫn giữa một pipeline
``tình cờ cho kết quả trông đúng'' trên dữ liệu đơn điệu với một kiến trúc xử lý
thực sự đúng đắn; tỷ lệ lỗi theo từng chỉ số trên cửa sổ thời gian chính là một
dạng quality meta-stream đơn giản hoá. Báo cáo không tích hợp framework Stream
DaQ đầy đủ (ngoài phạm vi bản một máy) mà vận dụng tinh thần của nó vào cơ chế
Quarantine và metric veracity sẵn có.

\textbf{Trạng thái tài liệu.} Stream DaQ~\cite{stream_daq_2024} được tham
chiếu theo preprint arXiv năm 2025, định danh \texttt{2506.06147}, DOI
\texttt{10.48550/arXiv.2506.06147} đã xác minh. Báo cáo không gán nơi
xuất bản tạp chí chưa có evidence; chỉ dùng khái niệm về phép kiểm tra và
quality meta-stream. Số liệu độ trễ hay thông lượng của công trình không
được dùng làm số đo của pipeline trong báo cáo này.

\subsection{Bảng so sánh khoảng trống nghiên cứu}
\label{subsec:gap-table}

Bảng~\ref{tab:gap} định vị báo cáo so với các công trình trên theo năm khả
năng: (i) xử lý event-time/watermark; (ii) chất lượng dữ liệu stream;
(iii) thích nghi runtime; (iv) chỉnh sửa/đối chiếu batch; (v) vận hành
IoT/lakehouse.

\begin{table}[H]
\centering
\caption{So sánh khoảng trống nghiên cứu giữa các công trình liên quan và báo cáo.
Ký hiệu: $\checkmark$ = có; một phần = có một phần; --- = không là trọng tâm.}
\label{tab:gap}
\footnotesize
\setlength{\tabcolsep}{3pt}
\begin{tabular}{p{4.6cm} p{1.9cm} p{1.9cm} p{1.9cm} p{1.9cm} p{1.9cm}}
\toprule
\textbf{Công trình} & \textbf{Event-time} & \textbf{Quality} & \textbf{Thích nghi} & \textbf{Correction} & \textbf{IoT/ lakehouse} \\
\midrule
Dataflow Model~\cite{dataflow_model_2015} & $\checkmark$ & --- & --- & --- & --- \\
MillWheel~\cite{millwheel_2013} & $\checkmark$ & --- & --- & --- & --- \\
Structured Streaming~\cite{structured_streaming_sigmod_2018} & $\checkmark$ & --- & --- & --- & --- \\
Strider~\cite{strider_iswc_2017} & một phần & --- & $\checkmark$ & --- & --- \\
Lambda ADS-B~\cite{lambda_adsb_2023} & --- & --- & --- & $\checkmark$ & một phần \\
Kafka patterns/benchmark~\cite{kafka_patterns_benchmark_2025} & --- & --- & --- & --- & một phần \\
Stream DaQ~\cite{stream_daq_2024} & --- & $\checkmark$ & một phần & --- & --- \\
\textbf{Báo cáo này} & một phần (baseline) & $\checkmark$ (veracity vector + per-metric) & đề xuất, chưa triển khai & một phần (batch Gold theo giờ) & $\checkmark$ \\
\bottomrule
\end{tabular}
\end{table}

Đọc bảng theo hàng (cột ``Event-time'' bao gồm watermark, ``Quality'' là chất
lượng dữ liệu stream, ``Thích nghi'' là thích nghi runtime, ``Correction'' là
chỉnh sửa/đối chiếu batch, ``IoT/lakehouse'' là vận hành trên miền
IoT/lakehouse): không có công trình đơn lẻ nào trong danh mục tham khảo
của báo cáo kết hợp đủ năm khả năng cho miền IoT. Nền tảng
event-time đã chín muồi~\cite{dataflow_model_2015, millwheel_2013} nhưng là mô
hình chung, không gắn với veracity; các công trình adaptive
~\cite{strider_iswc_2017} tập trung vào query plan của miền RDF stream; chỉnh
sửa batch~\cite{lambda_adsb_2023} gắn với kiến trúc Lambda đầy đủ. Hai công
trình bổ sung thu hẹp khoảng trống ở hai cột: thực hành benchmark cho hệ
event-streaming~\cite{kafka_patterns_benchmark_2025} cung cấp cơ sở phương pháp
đo tái lập được (danh mục công bố cấu hình tối thiểu, dạng đường cong
throughput theo partition/consumer), còn Stream DaQ~\cite{stream_daq_2024} cung
cấp mô hình giám sát chất lượng stream-first mà báo cáo vận dụng ở mức khái
niệm (phân loại phép kiểm tra, quality meta-stream) cho cơ chế Quarantine và
metric veracity per-metric. Khoảng trống còn lại mà báo cáo nhắm tới --- gắn
veracity vector đa chỉ số, metric rolling theo cửa sổ, policy thích nghi và
vòng đời công bố vào một pipeline MQTT--Kafka--Spark--Parquet/PostgreSQL cụ thể
cho cảm biến môi trường --- vẫn là khoảng trống hợp lệ về mặt định vị; tuy
nhiên báo cáo \emph{không} tuyên bố đã lấp đầy nó bằng triển khai hay số đo,
mới chỉ cung cấp định nghĩa, bất biến, giả thuyết kiểm chứng và một bản triển
khai tham chiếu trên một máy.

\textbf{Trạng thái nguồn tham khảo.} Metadata xuất bản của nghiên cứu
Kafka~\cite{kafka_patterns_benchmark_2025} đã được xác minh qua Crossref.
Stream DaQ~\cite{stream_daq_2024} được trích dẫn theo bản preprint arXiv
năm 2025, với định danh \texttt{2506.06147} và DOI arXiv đã xác minh; báo
cáo không gán cho tài liệu này một nơi xuất bản tạp chí chưa có evidence.
Phần tổng quan chỉ sử dụng mô hình chất lượng dữ liệu của công trình,
không chuyển số liệu hiệu năng của nó thành kết quả cho pipeline này.

Chương~\ref{chap:implementation} cụ thể hoá nền tảng này thành kiến trúc,
data contract và cấu hình vận hành; phần baseline khớp với mô hình
đã định nghĩa ở Mục~\ref{sec:theory}.
